

Comme l'anneau considéré est intègre, on en déduit la relation suivante, valable par substitution dans tout @@M0@@-anneau, et pour tous éléments @@M1@@ tels que @@M2@@ pour @@M3@@, @@M4@@ pour @@M5@@~;
@@M6@@


5.5. Soit maintenant @@M7@@ un @@M8@@-anneau quelconque, et @@M9@@ un élément de @@M10@@ tel que @@M11@@ pour @@M12@@. On va associer à @@M13@@ un nouveau @@M14@@-anneau @@M15@@.

On définit d'abord une nouvelle multiplication sur @@M16@@ en posant pour @@M17@@ :
@@M18@@

L'anneau @@M19@@ muni de cette multiplication n'est plus unitaire; on considère donc l'anneau @@M20@@, obtenu par adjonction d'une unité à @@M21@@. Ses opérations sont donc définies par :
@@M22@@
On considère @@M23@@ comme plongé dans cet anneau par @@M24@@. On définit alors une structure de pré-@@M25@@-anneau par :
@@M26@@
C'est une structure de pré-@@M27@@-anneau en vertu de (5.4.1). On notera @@M28@@ le pré-@@M29@@-anneau ainsi obtenu. Pour voir que c'est un @@M30@@-anneau, il suffit par substitution de faire la démonstration dans le cas où @@M31@@ est le @@M32@@-anneau engendré sur @@M33@@ par les générateurs @@M34@@, soumis aux relations @@M35@@ pour @@M36@@. On définit alors un homomorphisme @@M37@@, en
